\newpage
\section{FUTURE ENHANCEMENTS}

While the final prototype shows the viability of using RL for antenna alignment, there are some possible optimizations in hardware and software that may further increase precision and robustness.


\subsection{Mechanical and Hardware Improvements}

The current prototype is sufficient for concept demonstration, yet it is possible to make some adjustments to optimize performance. One possible issue in the current prototype setup is that there is a possibility for backlash in the pan axis and vibration after a step due to a loose coupling mechanism, poor quality of bearings, and a flexible antenna mount. It is recommended to use high-quality bearings, a better coupling, and a more rigid mounting for the antenna. Additionally, it is possible to improve the performance of the homemade slip ring by increasing the contact pressure and using a better conductive finish.

Using a higher-gain directional antenna and a calibrated RSSI front-end may help to increase the SNR of the measurements. Such improvement will be especially useful in indoor or obstructed areas where the system needs to distinguish between multiple local RSSI peaks.


\subsection{Algorithmic Improvements}

This project uses a small Q-learning controller because of its simplicity, interpretability, and ability to run on a low-cost microcontroller. Yet, some future work might focus on improving the learning process of the algorithm to cope with rapidly changing environment. Some possible optimization include:

\begin{itemize}
    \item \textbf{Online policy adaptation:} Allow limited online adjustment of Q-values or policy parameters in response to sudden environmental changes.
    \item \textbf{Function approximation:} Replace tabular policy representation with lightweight function approximation or a value compression to get a finer resolution without a dramatic increase in memory requirements.
    \item \textbf{Hybrid search strategy:} Combine a fast coarse search with RL to decrease the convergence time for weak signals.
    \item \textbf{Temporal state features:} Add some temporal state feature like an RSSI history, packet loss rate, or moving average slope to distinguish between temporary and stable improvements.
\end{itemize}


\subsection{Embedded Software Enhancements}

The firmware currently uses a deterministic blocking control loop, which was appropriate
for experimental validation. In future versions, the software can be improved by adopting
interrupt-assisted sampling or an RTOS-based scheduling model. This would enable parallel
handling of RSSI acquisition, telemetry logging, and actuator control, thereby improving
responsiveness without compromising measurement integrity.

Additional software enhancements may include automatic experiment logging to onboard
storage, wireless transfer of recorded trial data, remote parameter tuning, and a simple
user interface for selecting alignment modes and calibration routines.

\subsection{System-Level Extensions}

The proposed system can be extended beyond the present laboratory-scale implementation.
A multi-node version could coordinate several receivers or repeaters simultaneously,
allowing distributed beam steering in wireless mesh or relay networks. The alignment logic
could also be integrated with link-layer performance indicators such as throughput, packet
error rate, or signal-to-noise ratio, thereby enabling the controller to optimize overall
communication quality rather than RSSI alone.

\subsection{Enhancements to the Embedded Software}

Currently, the embedded firmware utilizes a deterministic blocking control loop that was proven to work well during experimental evaluations. The next versions of the software might benefit from either adding support for interrupt-driven RSSI sampling or implementing the RTOS scheduling model that would allow RSSI reception, telemetry recording, and actuator movement to be performed simultaneously.

Additional possible enhancements of the firmware include automated experiment logging to the internal storage device, wireless uploading of collected experiments, online parameter tuning capabilities, and a simple user interface for choosing the alignment algorithm.

\subsection{Extensions to the System}

The described system architecture is not limited to its present lab-scale implementation only. It can be scaled up and used to implement multi-node systems where multiple receivers or repeaters are operated at once to provide beam steering in a mesh or relay wireless network. In addition, the alignment logic could be linked to some performance measures, such as throughput, packet loss rate, or SNR, thus optimizing the overall link performance, rather than just RSSI values.